% sections/ndss_09_conclusion.tex  — NDSS 2027

We presented \sysname, the first large-scale measurement of
capability-description inconsistency in LLM agent skill marketplaces,
and \cape, an empirical validation framework demonstrating
the real-world exploitability of the identified gaps.

Our analysis of \NskillsCH \clawhub skills reveals that \ICone
completely omit security-sensitive capabilities---shell execution,
credential access, file deletion, and data exfiltration---from visible
descriptions, yielding a mean Capability Recall of \MeanCR{}
(Cohen's $\kappa = \IAAsink$ inter-annotator agreement on ground truth).
This transparency failure is not latent risk.
\cape validates all \NpairsTotal source-sink pairs against real OpenClaw agents.
Under the pre-defined intent-to-treat metric (counting \NTimeout{} infrastructure
timeouts as failures), \NpairsSuccessITT/\NpairsTotal = \PairsSuccessRateITT{} succeed.
Under the completed-run metric (conditioned on agent returning a response),
\NpairsSuccessCR/\NpairsCompletedRun = \textbf{\PairsSuccessRateCR}
(95\% CI: \CICRlo--\CICRhi), with shell-execution sinks exploited at
\ShellSuccessRate{} (\NShellCompleted/\NShellCompleted{} completed pairs).
Only \NRefusal{} pairs (3.2\%) produce explicit agent refusals,
and these refusals are stochastic---the same task type succeeds in
67--100\% of other pairs---indicating the absence of a consistent safety policy.

Three structural findings stand out beyond the aggregate numbers.
First, the agent's safety policy is heuristic-based and inconsistently applied:
shell command execution via injected context is completely unguarded,
while direct memory-modification requests are occasionally refused
without a stable policy threshold.
Second, all three tested payload families (P1: \texttt{go vet}, P2: \texttt{node verify.js},
P3: \texttt{cat README}) produce comparable success rates,
confirming that the vulnerability is in the agent's trust model,
not in any property of a specific command syntax.
Third, the 64 infrastructure timeouts have a single root cause
(LLM context compaction at high prompt token usage)---
a measurement challenge, not an agent defense---and are being retried.

LLM skill marketplaces are operating on an honor system for
capability disclosure. As agent autonomy expands, this gap produces
compounding risk. We call on platform operators to adopt
machine-readable capability manifests, publish-time screening,
and consistent capability-based runtime policies as a minimum response.
Our tooling and dataset are released to facilitate this work.

\paragraph{Acknowledgments.}
Omitted for anonymous review.
